% Options for packages loaded elsewhere
\PassOptionsToPackage{unicode}{hyperref}
\PassOptionsToPackage{hyphens}{url}
%
\documentclass[
]{article}
\usepackage{amsmath,amssymb}
\usepackage{iftex}
\ifPDFTeX
  \usepackage[T1]{fontenc}
  \usepackage[utf8]{inputenc}
  \usepackage{textcomp} % provide euro and other symbols
\else % if luatex or xetex
  \usepackage{unicode-math} % this also loads fontspec
  \defaultfontfeatures{Scale=MatchLowercase}
  \defaultfontfeatures[\rmfamily]{Ligatures=TeX,Scale=1}
\fi
\usepackage{lmodern}
\ifPDFTeX\else
  % xetex/luatex font selection
\fi
% Use upquote if available, for straight quotes in verbatim environments
\IfFileExists{upquote.sty}{\usepackage{upquote}}{}
\IfFileExists{microtype.sty}{% use microtype if available
  \usepackage[]{microtype}
  \UseMicrotypeSet[protrusion]{basicmath} % disable protrusion for tt fonts
}{}
\makeatletter
\@ifundefined{KOMAClassName}{% if non-KOMA class
  \IfFileExists{parskip.sty}{%
    \usepackage{parskip}
  }{% else
    \setlength{\parindent}{0pt}
    \setlength{\parskip}{6pt plus 2pt minus 1pt}}
}{% if KOMA class
  \KOMAoptions{parskip=half}}
\makeatother
\usepackage{xcolor}
\usepackage[margin=1in]{geometry}
\usepackage{graphicx}
\makeatletter
\def\maxwidth{\ifdim\Gin@nat@width>\linewidth\linewidth\else\Gin@nat@width\fi}
\def\maxheight{\ifdim\Gin@nat@height>\textheight\textheight\else\Gin@nat@height\fi}
\makeatother
% Scale images if necessary, so that they will not overflow the page
% margins by default, and it is still possible to overwrite the defaults
% using explicit options in \includegraphics[width, height, ...]{}
\setkeys{Gin}{width=\maxwidth,height=\maxheight,keepaspectratio}
% Set default figure placement to htbp
\makeatletter
\def\fps@figure{htbp}
\makeatother
\setlength{\emergencystretch}{3em} % prevent overfull lines
\providecommand{\tightlist}{%
  \setlength{\itemsep}{0pt}\setlength{\parskip}{0pt}}
\setcounter{secnumdepth}{-\maxdimen} % remove section numbering
\usepackage{wasysym}
\ifLuaTeX
  \usepackage{selnolig}  % disable illegal ligatures
\fi
\usepackage{bookmark}
\IfFileExists{xurl.sty}{\usepackage{xurl}}{} % add URL line breaks if available
\urlstyle{same}
\hypersetup{
  pdftitle={Mathematical description of the IBM},
  hidelinks,
  pdfcreator={LaTeX via pandoc}}

\title{Mathematical description of the IBM}
\author{}
\date{\vspace{-2.5em}2025-05-09}

\begin{document}
\maketitle

We developed a discrete-time, discrete-space stochastic individual-based
model (IBM) to simulate the population dynamics, genetic structure
(including inheritance), and dispersal of a generic species. The model
accommodates both overlapping and non-overlapping generations. We use
the model to investigate whether genetic load as a result of the
accumulation of deleterious alleles in the population can generate
generate Allee effect to slow down the invasion process. We then explore
the impact of different spatial topologies (one-dimensional space using
a stepping-stone model vs a two-dimensional metapopulation model. The
stepping-stone model represents a connection of linear patches where
dispersal of individuals between patches is limited to nearest
neighbours(adjacent patches). In contrast, the metapopulation model is a
network of patches connected by dispersal. Both models incorporate key
biological processes: mating/reproduction, survival, and dispersal. We
assume all patches are homogeneous in terms of suitability. Simulations
are undertaken on a yearly time step indexed by \(t\). Within each patch
indexed by \(j\) are individuals indexed by \(i\). For non-overlapping
generations, individuals reproduce and die after each timestep, in
contrast to overlapping generations.

\section{Reproduction}\label{reproduction}

We modelled a hermaphroditic panmictic population where all individuals
are capable of reproducing by breeding with a randomly selected mate in
the population. Each individual has an average reproductive output,
\(f\) (fecundity). We modelled reproduction in a density-dependent
manner using the Beverton--Holt model (Beverton \& Holt 2012), which
yields an expected reproductive output, \(E(f)\). The actual
reproductive output \(A(f)\) for the population is then drawn from a
Poisson distribution.

\[
A(f)\sim \text{Poisson}(N,E(f))
\]

where: \[
E(f) = \frac{f}{1 + (\frac{f - 1}{K}) \times N}
\]

where \(K\) is the carrying capacity and \(N\) is the population size.

\section{Survival}\label{survival}

We assume that survival rate is the same for every individual in the
population and that they survive with a probability \(s\). Thus, whether
an individual survives or not \(S_{i,j,t}\) at each time step \(t\) is
modelled as a stochastic process and drawn from a Bernoulli
distribution. However, in the non-overlapping generation, following
reproduction, all parents die with a 100\% probability.

\[
S_{i,j,t} \sim\text{Bernoulli} (s).
\]

\section{Dispersal}\label{dispersal}

Dispersal is a fundamental driver of invasion dynamics. As individuals
seek resources, mates, or evade predators, they expand their range,
consequently facilitating the spread and establishment in new areas. We
model dispersal using two approaches: a negative exponential dispersal
kernel for the meta-population model and an adjacency matrix
(nearest-neighbour) for the stepping-stone model. During dispersal,
dispersing individuals move from an origin \(j\) to a destination \(k\)
patch, with a probability. In the metapopulation model, this probability
decreases exponentially as the distance between patches increases. In
contrast, the stepping-stone model assumes that individuals disperse to
adjacent patches with equal probability. The dispersal kernel \(K\) is
given below as:

\[
K_{j,k} = 
\begin{cases}
dx, & j \ne k \\
0, & j = k
\end{cases}
\] Where \(dx\) can either denote a negative exponential dispersal
matrix \(e^{-\lambda W_{j,k}}\), or an adjacency matrix \(U_{j,k}\).
Each row of the resultant dispersal matrix is then normalised to produce
conditional probabilities \(O_{j,k}\):

\[
O_{j,k} = \frac{dx} {\sum_{k} dx}
\] and given each individual has a daily probability \(r\) of
dispersing, the true probability \(X_{j,k}\) of moving from \(j\) to
\(k\) is calculated as:

\[
X_{j,k} =
\begin{cases}
rO_{j,k}, & j \ne k \\
1-r, & j = k 
\end{cases}
\]

\section{Genetic inheritance}\label{genetic-inheritance}

The gametes for each offspring are formed by selecting one allele per
locus from each parent. Rather than assuming independent segregation, we
modelled genetic linkage by generating correlated allele choices. To do
this, we obtained the distance matrix \(Q_{x,y}\) between pairs of loci
based on their Euclidean distances and calculated the
variance-covariance matrix \(\Sigma\) using a covariance decay function.

\[
\Sigma = \sigma^2 \exp(−λQ_{x,y})
\]

where \(\sigma^2\) is a scaling factor and \(\lambda\) controls the
decay rate of the covariance between two loci as genetic distance
increases. Using \(\Sigma\), we then generated multivariate normal
random effect values \(\epsilon\) that were transformed to probabilities
via a logistic function. The resulting probabilities were then used to
conduct Bernoulli trials to determine which allele (from each parent) is
inherited by the offspring for each locus. Thus, if \(L\) is the number
of loci, and \(\Sigma\) the variance-covariance matrix describing
correlations among loci, the multivariate random effect \(\epsilon\) is
given as:

\[
\epsilon_i \sim \mathcal{N}_L\left(\mu, \Sigma \right)
\]

where \(\epsilon_i = (\epsilon_{i1}, \dots, \epsilon_{iL})\) and
\(\mathcal{N}_L\) denotes the multivariate normal distribution, and
\(\mu\) is a vector with \(L\) dimensional mean of 0. Consequently, the
probability \(\rho_{i}\) of inheriting an allele in any locus from a
given parent using the logistic function becomes, where
\(\rho_{i}  = (\rho_{il} , \dots,\rho_{iL})\).

\[
\rho_{i} = \frac{1}{1 + (\exp (-\epsilon_{i}))}
\]

\end{document}
